\documentclass[11pt]{article}
\usepackage[margin=1in]{geometry}
\usepackage{amsmath,amssymb,amsthm}
\usepackage{hyperref}

\newtheorem{proposition}{Proposition}
\newcommand{\blank}{\varnothing}
\newcommand{\Py}{P_{\mathrm{y}}}
\newcommand{\Pc}{P_{\mathrm{c}}}
\newcommand{\Peff}{P_{\mathrm{eff}}}

\title{A Confidence Head in the RNNT Forward Recursion}
\author{Working note}
\date{\today}

\begin{document}
\maketitle

\section{Motivation}

In a full-context transducer, $P(\text{token} \mid t, \text{history})$ is well posed:
the encoder has observed the whole utterance, so the time index $t$ merely selects
a position in a fully observed sequence. The blank symbol $\blank$ then carries a
single, unambiguous meaning: \emph{nothing is emitted here, advance time}.

In the streaming setting the encoder at chunk $t$ has observed audio only up to
chunk $t$. The blank symbol becomes overloaded, conflating two situations that
call for different downstream behaviour:

\begin{enumerate}
  \item \textbf{Confident emptiness.} There is genuinely nothing to emit at this
        chunk. The decision is final and will not be revised.
  \item \textbf{Deferral under uncertainty.} The model cannot yet determine what
        to emit and would prefer to observe more audio before committing.
\end{enumerate}

Both produce the same symbol, the same lattice transition, and the same gradient.
The proposal examined here is to add a second output head emitting an explicit
binary confidence at every lattice position, so that the two cases become
distinguishable.

\section{Notation}

Let $X$ be the input and $Y = (y_1,\dots,y_U)$ the target token sequence. The
lattice node $(t,u)$ denotes ``$t$ chunks consumed, $u$ tokens emitted'', with
$1 \le t \le T$ and $0 \le u \le U$.

\paragraph{Token head.} $\Py(k \mid t,u)$ is a distribution over
$V \cup \{\blank\}$, where $V$ is the token vocabulary.

\paragraph{Confidence head.} $\Pc(1 \mid t,u)$ and $\Pc(0 \mid t,u)$ with
$\Pc(1)+\Pc(0)=1$, interpreted as confident / not confident.

Where no ambiguity arises the arguments $(t,u)$ are suppressed.

\section{The standard recursion}

With a single head, the forward variable
$\alpha(t,u) = P(\text{reach } (t,u))$ obeys
\begin{equation}
  \alpha(t,u) \;=\; \alpha(t-1,u)\,\Py(\blank \mid t-1,u)
                 \;+\; \alpha(t,u-1)\,\Py(y_u \mid t,u-1),
  \label{eq:rnnt}
\end{equation}
with $\alpha(1,0)=1$, and the sequence likelihood
$P(Y \mid X) = \alpha(T,U)\,\Py(\blank \mid T,U)$.

Every path through the lattice is one alignment; \eqref{eq:rnnt} sums over all of
them. Crucially, in the full-sum formulation the alignment is \emph{latent}: no
supervision identifies which chunk a given word ``should'' have been emitted at.

\section{First formulation: confidence gates commitment}
\label{sec:gating}

The natural reading of the confidence bit is that it gates committing. If
confident, the model emits a token or a genuine blank; if not confident, it
declines to emit and advances time. The transitions out of $(t,u)$ are then

\begin{center}
\begin{tabular}{lll}
\textbf{action} & \textbf{probability} & \textbf{successor} \\[2pt]
emit $y_{u+1}$                 & $\Pc(1)\,\Py(y_{u+1})$ & $(t,\,u+1)$ \\
blank, confident               & $\Pc(1)\,\Py(\blank)$  & $(t+1,\,u)$ \\
blank, not confident (defer)   & $\Pc(0)$               & $(t+1,\,u)$ \\
\end{tabular}
\end{center}

\noindent giving the recursion
\begin{equation}
  \alpha(t,u) \;=\; \alpha(t-1,u)\big[\Pc(1)\Py(\blank) + \Pc(0)\big]
                 \;+\; \alpha(t,u-1)\,\Pc(1)\Py(y_u).
  \label{eq:gated}
\end{equation}

\subsection{This formulation is unidentifiable}

\begin{proposition}
Under \eqref{eq:gated} the likelihood $P(Y \mid X)$ depends on the two heads only
through the composite quantities
\begin{equation}
  \Peff(\blank) \;=\; \Pc(1)\Py(\blank) + \Pc(0),
  \qquad
  \Peff(y) \;=\; \Pc(1)\Py(y) \quad (y \in V).
  \label{eq:peff}
\end{equation}
Consequently infinitely many pairs $(\Pc,\Py)$ attain the same loss, and the
confidence head is not identified.
\end{proposition}

\begin{proof}
First, $\Peff$ is a valid distribution:
\[
  \sum_{y \in V}\Peff(y) + \Peff(\blank)
  = \Pc(1)\Big[\sum_{y\in V}\Py(y) + \Py(\blank)\Big] + \Pc(0)
  = \Pc(1) + \Pc(0) = 1 .
\]
Substituting \eqref{eq:peff} into \eqref{eq:gated} yields
\[
  \alpha(t,u) = \alpha(t-1,u)\,\Peff(\blank \mid t-1,u)
              + \alpha(t,u-1)\,\Peff(y_u \mid t,u-1),
\]
which is exactly \eqref{eq:rnnt} with $\Py$ replaced by $\Peff$. Hence the loss is
a function of $\Peff$ alone. The defining relation
$\Peff(\blank) = \Pc(1)\Py(\blank) + \Pc(0)$ is a single equation in the two
unknowns $\Pc(0)$ and $\Py(\blank)$, so for any target $\Peff(\blank) \in (0,1)$
there is a one-parameter family of solutions, all with identical likelihood.
\end{proof}

\paragraph{Consequences.}
\begin{itemize}
  \item Gradient descent places $\Pc$ arbitrarily within the family; whatever it
        converges to reflects initialisation and optimiser dynamics, not a
        property of the data.
  \item Adding a cost $\lambda\,\Pc(0)$ to discourage deferral does not repair
        this. Both blank transitions reach the same successor $(t+1,u)$, so the
        deferral branch confers no compensating benefit and the penalised
        optimum is simply $\Pc(0) \equiv 0$.
  \item The underlying reason is structural: \emph{two actions that lead to the
        same successor state and differ only in label are exchangeable.} Any fix
        must make the two blanks differ in consequence, not merely in name.
\end{itemize}

\section{Second formulation: deferral with a budget}
\label{sec:budget}

Augment the lattice state with a counter $d$ recording the number of
\emph{consecutive} deferrals, $0 \le d \le k$ for a fixed budget $k$. The node is
now $(t,u,d)$. Transitions out of $(t,u,d)$:

\begin{center}
\begin{tabular}{llll}
\textbf{action} & \textbf{probability} & \textbf{successor} & \textbf{condition} \\[2pt]
emit $y_{u+1}$               & $\Pc(1)\Py(y_{u+1})$ & $(t,\,u+1,\,0)$ & --- \\
blank, confident             & $\Pc(1)\Py(\blank)$  & $(t+1,\,u,\,0)$ & --- \\
defer                        & $\Pc(0)$             & $(t+1,\,u,\,d+1)$ & $d < k$ \\
\end{tabular}
\end{center}

At $d = k$ the deferral branch is removed and the confidence head is masked so
that $\Pc(1) \equiv 1$; the distribution therefore remains normalised at every
node. Both committing actions reset the counter.

The forward recursion splits by whether the counter is zero:
\begin{align}
  \alpha(t,u,0) &= \sum_{d'=0}^{k} \alpha(t-1,u,d')\,
                     \Pc(1 \mid t-1,u,d')\,\Py(\blank \mid t-1,u,d') \notag \\
                &\quad + \sum_{d'=0}^{k} \alpha(t,u-1,d')\,
                     \Pc(1 \mid t,u-1,d')\,\Py(y_u \mid t,u-1,d'),
  \label{eq:budget0} \\[4pt]
  \alpha(t,u,d) &= \alpha(t-1,u,d-1)\,\Pc(0 \mid t-1,u,d-1),
                   \qquad 1 \le d \le k,
  \label{eq:budgetd}
\end{align}
with $\alpha(1,0,0)=1$ and all other $\alpha(1,\cdot,\cdot)=0$. The likelihood is
\begin{equation}
  P(Y \mid X) \;=\; \sum_{d=0}^{k} \alpha(T,U,d)\,
                    \Pc(1 \mid T,U,d)\,\Py(\blank \mid T,U,d),
\end{equation}
requiring the sequence to terminate on a confident blank.

\subsection{Why this is identifiable}

Deferral now has a consequence: it consumes budget, and budget exhaustion
forecloses continuations that would otherwise be available. Two parameterisations
that agree on $\Peff$ but differ in how mass splits between $\Pc(0)$ and
$\Py(\blank)$ reach different distributions over $d$, hence different sets of
reachable paths, hence different likelihoods. The collapse of
Section~\ref{sec:gating} no longer applies.

The limiting behaviour is a useful check. As $k \to \infty$ the budget never
binds, the counter becomes irrelevant, and the model degenerates to
\eqref{eq:gated} --- unidentifiable again. \emph{Finiteness of $k$ is exactly what
makes the confidence head learnable.}

\subsection{Cost}

The lattice grows by a factor of $k+1$. For the intended regime ($k = 2$ or $3$)
this is a small constant. Equations \eqref{eq:budget0}--\eqref{eq:budgetd} retain
the same dependency structure as \eqref{eq:rnnt}, so the usual parallel
forward--backward schedule over anti-diagonals still applies.

\subsection{Interpretation of $k$}

$k$ is simultaneously (i) the regulariser that makes $\Pc$ identifiable and
(ii) a hard upper bound on emission latency, in chunks. At chunk duration
$0.08\,\mathrm{s} \times 14 = 1.12\,\mathrm{s}$, a budget of $k=2$ bounds extra
deferral at $2.24\,\mathrm{s}$.

\section{Open questions}

\begin{enumerate}
  \item \textbf{Does this merely relabel the band?} A budget of $k$ permits a word
        to be emitted at most $k$ chunks after it first could have been, which is
        close to what a fixed alignment band enforces --- except allocated by the
        model rather than fixed in advance. The claim worth testing is that
        \emph{learned, input-dependent} allocation beats a \emph{fixed} band at
        equal mean latency. This is directly testable against an existing banded
        arm.

  \item \textbf{Does $\Pc$ mean what we want it to mean?} Nothing in the objective
        forces $\Pc(0)$ to track uncertainty; it marks ``took the budgeted path''.
        Whether that coincides with low per-token correctness is empirical, and
        should be measured (e.g.\ correlation of $\Pc$ with per-word correctness,
        and calibration, not merely ranking) rather than assumed.

  \item \textbf{Where is the confidence head conditioned?} The recursions above
        assume $\Pc$ is a function of the same node representation as $\Py$. If
        $\Pc$ is instead conditioned on a narrower or wider context, the
        identifiability argument is unchanged but the inductive bias differs.

  \item \textbf{Inference.} The threshold on $\Pc$ is an inference-time concern and
        is best kept out of the trained model, so that operating points can be
        swept on a fixed checkpoint. Evaluation should be a WER / mean-latency
        curve rather than a single operating point.
\end{enumerate}

\end{document}
